% LaTeX version of PLAN.md, for sharing with members who have not joined the GitHub org yet.
% PLAN.md is the source of truth: when it changes, update this file to match.
% Build: latexmk -pdf plan.tex   (or: pdflatex plan.tex, twice)
\documentclass[11pt,a4paper]{article}

\usepackage[T1]{fontenc}
\usepackage[utf8]{inputenc}
\usepackage{lmodern}
\usepackage{microtype}
\usepackage{upquote}
\usepackage[margin=2.2cm]{geometry}
\usepackage{booktabs}
\usepackage{tabularx}
\usepackage{array}
\usepackage{enumitem}
\usepackage[table]{xcolor}
\usepackage{tcolorbox}
\usepackage{titlesec}
\usepackage{fancyhdr}
\usepackage[hidelinks]{hyperref}

\definecolor{accent}{HTML}{1F4E79}
\hypersetup{colorlinks=true,linkcolor=accent,urlcolor=accent}
\titleformat{\section}{\Large\bfseries\color{accent}}{\thesection.}{0.5em}{}
\titleformat{\subsection}{\large\bfseries}{}{0em}{}
\setlist{itemsep=2pt,topsep=4pt}
\setlength{\parindent}{0pt}
\setlength{\parskip}{6pt}
\renewcommand{\arraystretch}{1.25}
\newcolumntype{L}{>{\raggedright\arraybackslash}X}

\pagestyle{fancy}
\fancyhf{}
\lhead{\small IS 201 Final Project Plan}
\rhead{\small Ctrl+Alt+Ethics}
\cfoot{\small\thepage}
\renewcommand{\headrulewidth}{0.4pt}

\newcommand{\code}[1]{\texttt{#1}}
\newcommand{\repo}{https://github.com/is-201-ctrl-alt-ethics/genai-ethics-simulator}

\begin{document}

\begin{center}
  {\LARGE\bfseries\color{accent} GenAI Ethics Simulator: Project Plan}\par\vspace{4pt}
  {\large IS 201: Computer Ethics \textperiodcentered{} Final Group Project}\par\vspace{2pt}
  Ctrl+Alt+Ethics \textperiodcentered{} University of the Philippines \textperiodcentered{} First Semester, AY 2026--2027
\end{center}

\begin{tcolorbox}[colback=accent!6,colframe=accent,boxrule=0.6pt,arc=2pt]
Working plan for a \textbf{Technology and Society Ethics Simulator} for
\textbf{Scenario 4: Generative AI in Universities} (institutional policy on GenAI in teaching,
assessment and administration).

\textbf{Due: Saturday, December 19, 2026.} Submit a working website link and a 10--12 minute YouTube video.

\begin{tabular}{@{}ll@{}}
Repository: & \url{\repo} \\
Live site:  & \url{https://is-201-ctrl-alt-ethics.github.io/genai-ethics-simulator/} \\
Full brief: & \code{docs/requirements/final-req.md} in the repository \\
\end{tabular}
\end{tcolorbox}

\section{What the requirements boil down to}

\begin{tabularx}{\textwidth}{@{}p{3.1cm}LL@{}}
\toprule
\textbf{Area} & \textbf{Hard requirement} & \textbf{Where it lives in the repo} \\
\midrule
Decision variables & 5+, each with 2+ meaningful options & \code{data/model.json} $\rightarrow$ \code{variables} \\
Indicators & 5+, distinct, clear interpretation; must cover \textbf{equity}, \textbf{governance}, \textbf{sustainability} + others & \code{data/model.json} $\rightarrow$ \code{indicators} \\
Interactions & 4+, each on a \textbf{different pair} of variables, each justified & \code{data/model.json} $\rightarrow$ \code{interactions} \\
Course concepts & The 3 central areas \textbf{shape the model}, plus 2+ of: ethical frameworks and professional responsibility, privacy and data rights, cybersecurity, AI, software development & Tagged per rule in \code{model.json}; explained on the site \\
Model behaviour & Auto-updates; several variables per indicator; real trade-offs (no ``right answer''); understandable rules & \code{js/engine.js} + the rendered rules table \\
Site sections & Scenario overview + stakeholders; decision interface; dashboard; ethical explanations; model transparency (rules, assumptions, limitations, hypothetical labels); references & \code{index.html} sections \\
Evidence & Each key relationship: \textbf{why}, \textbf{what supports it}, \textbf{assumptions and limitations}; how weights were set & \code{justification} + \code{sources} per rule; \code{data/references.json} \\
No separate report & Everything is inside the site & Nothing graded lives only in \code{docs/} \\
Video & 10--12 min; scenario, model + justification, demo, insights, limitations; \textbf{every member speaks} & \code{docs/video/script.md} \\
\bottomrule
\end{tabularx}

\textbf{Grading:} ethical reasoning 30\% \textperiodcentered{} model quality and justification 30\%
\textperiodcentered{} functionality and usability 25\% \textperiodcentered{} presentation 15\%.
So 60\% rides on the \emph{reasoning and evidence}, not the code, and the plan front-loads research
and model design.

\section{Key design decision: one model file drives everything}

The model (variables, options, indicators, base values, effects, interactions, justifications and
source IDs) lives in \textbf{one data file}, \code{data/model.json}. The site:
\begin{enumerate}
  \item builds the decision controls from it,
  \item computes the indicators from it (\code{js/engine.js}, a small pure function), and
  \item renders the ``How the model works'' rules table and the justifications from it.
\end{enumerate}

So the rules users read are \emph{exactly} the rules that run, and the team can change the model
without touching UI code. A small check script (\code{tools/check-model.js}, run in CI) fails if the
model drops below the minimums: under 5 variables or indicators, a variable with fewer than 2
options, a missing required dimension, fewer than 4 interactions, two interactions on the same
pair, or a rule with no justification or source.

Scoring follows the brief's sample: every indicator starts at a base (e.g.\ 50), rule effects add
up, and the result is clamped to 0--100. All values are labelled \textbf{hypothetical}.

\section{Repo structure}

\begin{verbatim}
genai-ethics-simulator/
|-- index.html              # the whole site, served by GitHub Pages
|-- css/style.css
|-- js/
|   |-- engine.js           # pure: (model, decisions) -> scores + rules fired
|   `-- app.js              # controls, dashboard, explanations, rules table
|-- data/
|   |-- model.json          # single source of truth for the model
|   `-- references.json     # citation list; rules cite these by id
|-- tools/check-model.js    # requirement checks on model.json (run by CI)
|-- docs/                   # team working files, NOT graded
|   |-- requirements/final-req.md
|   |-- research/           # evidence log: one note per source or topic
|   `-- video/script.md     # video outline, speaker split
|-- .github/workflows/check.yml
|-- PLAN.md
`-- README.md
\end{verbatim}

Site sections in \code{index.html}: \textbf{Overview} (context, stakeholders) $\rightarrow$
\textbf{Simulator} (decisions, live dashboard, which rules fired) $\rightarrow$ \textbf{Trade-offs}
(ethical explanations) $\rightarrow$ \textbf{How the model works} (rules table, assumptions,
limitations, how weights were set) $\rightarrow$ \textbf{References}.

\section{Draft model (a starting point for the team, not a decision)}

These are \emph{candidates} to argue over in the model-design meeting. Every item kept needs
evidence (Section~\ref{sec:evidence}).

\subsection{Stakeholders}
Students (including those with limited devices or connectivity, and non-native English writers),
faculty, administrators and staff, the university as an institution, AI vendors, parents and the
public, and the environment.

\subsection{Candidate decision variables}
\begin{tabularx}{\textwidth}{@{}p{3.6cm}L@{}}
\toprule
\textbf{Variable} & \textbf{Options} \\
\midrule
Usage policy & Ban \textperiodcentered{} Allowed with disclosure \textperiodcentered{} Fully integrated \\
AI detection & None \textperiodcentered{} Detector as a flag for human review \textperiodcentered{} Detector output as evidence \\
Tool provisioning & Students bring their own (free tiers) \textperiodcentered{} Institution-licensed enterprise tool \textperiodcentered{} Self-hosted open model \\
Assessment design & Unchanged \textperiodcentered{} Partly redesigned (oral, in-class, process-based) \textperiodcentered{} Fully redesigned \\
Data governance & Vendor default terms \textperiodcentered{} Negotiated data agreement (no training on student data) \textperiodcentered{} On-premise only \\
AI literacy training & None \textperiodcentered{} Optional \textperiodcentered{} Mandatory for students and faculty \\
Administrative use & None \textperiodcentered{} Drafting and support only \textperiodcentered{} Automated decisions (e.g.\ admissions or advising triage) \\
Policy governance & Administration only \textperiodcentered{} Committee with faculty and student representatives \\
\bottomrule
\end{tabularx}

\subsection{Candidate indicators}
\begin{tabularx}{\textwidth}{@{}p{6cm}L@{}}
\toprule
\textbf{Indicator} & \textbf{Dimension} \\
\midrule
Equity of access & Equity and social justice \\
Fairness of integrity enforcement & Equity and social justice \\
Accountability and transparency & Technology governance \\
Environmental footprint & Computing and sustainability \\
Student data privacy & Privacy, data and digital rights \\
Learning and skill development & Other (learning outcomes) \\
Assessment validity & Other (academic integrity) \\
Faculty workload & Other (professional responsibility) \\
\bottomrule
\end{tabularx}

\subsection{Candidate interactions (each on a different pair)}
\begin{enumerate}
  \item \textbf{AI detection $\times$ Usage policy:} detectors under a ban drive more accusations,
        and known false positives against non-native English writers hurt \emph{fairness of
        enforcement} more.
  \item \textbf{Tool provisioning $\times$ Usage policy:} full integration with student-supplied
        tools widens the paid-tier gap (\emph{equity}); with institution-licensed tools it narrows it.
  \item \textbf{Tool provisioning $\times$ Data governance:} negotiated terms protect
        \emph{privacy} only for institution-licensed tools; with bring-your-own tools the university
        controls nothing.
  \item \textbf{AI literacy $\times$ Assessment design:} redesigned assessment raises
        \emph{learning} more when students and faculty are trained to use the tools well.
  \item \textbf{Administrative use $\times$ Policy governance:} automated admin decisions hurt
        \emph{accountability} far more without a representative oversight body.
  \item A \emph{sustainability} interaction: \textbf{Tool provisioning $\times$ Usage policy} on
        energy use is taken by \#2, so look for another pair (e.g.\ self-hosted hardware $\times$
        usage level). Keep every pair distinct.
\end{enumerate}

\subsection{Course concepts to tag}
Central: equity, governance, sustainability. Plus \textbf{AI}, \textbf{privacy, data and digital
rights}, and \textbf{ethical frameworks and professional responsibility} (e.g.\ ACM Code of Ethics,
faculty duty of care).

\section{Evidence to find and verify}\label{sec:evidence}

A member checks every source before it is cited. Leads to start from (not yet verified):
\begin{itemize}
  \item UNESCO, \emph{Guidance for generative AI in education and research} (2023).
  \item Liang et al.\ (2023), GPT detectors are biased against non-native English writers (\emph{Patterns}).
  \item Republic Act 10173, the Data Privacy Act of 2012 (Philippines), and National Privacy Commission issuances.
  \item University of the Philippines guidance on responsible AI use, and any CHED guidance (check what exists).
  \item EU AI Act, which lists some education uses as high-risk (a governance comparison).
  \item Research on the energy and water use of AI training and inference (e.g.\ Luccioni et al.).
  \item ACM Code of Ethics and Professional Conduct.
  \item Documented cases: universities that banned and later reversed GenAI policies; false
        accusations from AI detectors.
\end{itemize}
Log each source in \code{docs/research/} with the claim it supports, the rule(s) it backs, and its
limitations.

\section{Timeline (milestones)}

\begin{tabularx}{\textwidth}{@{}>{\raggedright\arraybackslash}p{4.4cm}p{2.7cm}L@{}}
\toprule
\textbf{Milestone} & \textbf{Dates} & \textbf{Outcome} \\
\midrule
\textbf{M1} Research and model design & Nov 24 -- Dec 1 & Stakeholders, final variables, indicators and interactions, evidence log, weights rationale \\
\textbf{M2} Simulator build & Dec 2 -- Dec 8 & \code{model.json} filled in; engine, UI, dashboard and rules table working on Pages; CI check passing \\
\textbf{M3} Site content & Dec 5 -- Dec 12 & Overview, trade-off explanations, assumptions, limitations and references written into the site \\
\textbf{M4} Review and polish & Dec 12 -- Dec 15 & Full checklist audit, usability test with someone outside the group, mobile check, fixes \\
\textbf{M5} Video and submission & Dec 15 -- Dec 19 & Script, recording (everyone speaks), YouTube upload, both links submitted \\
\bottomrule
\end{tabularx}

Aim to submit by \textbf{December 18} to keep a day of buffer.

\section{How we work}
\begin{itemize}
  \item Work is tracked as issues under the milestones above, using the org's issue forms (Task,
        Bug, Feature, Docs): \url{\repo/issues}.
  \item Branch per issue, open a pull request, and get one teammate's review before merging to
        \code{main}. \code{main} is what GitHub Pages serves.
  \item \textbf{GenAI use:} allowed for research help, coding and debugging, per the brief. Members
        verify every source, and the model's rules and ethical judgments are the group's own; AI
        output never stands in for a justification.
  \item Every member contributes and appears in the video, as the brief requires.
\end{itemize}

\section{Open questions}
\begin{itemize}
  \item Team members' GitHub handles and roles (research, model, front-end, content, video).
  \item Final variable and indicator set, and the weighting scheme (fixed points as in the sample,
        or something else).
  \item One page or several? One page keeps the demo simple, so that is the default.
\end{itemize}

\end{document}
